\documentclass[11pt,a4paper]{article}

% ---- Encoding and Fonts ----
\usepackage[UTF8]{ctex}
\setCJKmainfont{SimSun}         % 正文宋体
\setCJKsansfont{SimHei}         % 无衬线黑体
\setCJKmonofont{FangSong}       % 等宽仿宋
\usepackage[T1]{fontenc}
\usepackage{lmodern}

% ---- Layout ----
\usepackage[margin=2.5cm]{geometry}
\usepackage{setspace}
\onehalfspacing

% ---- Packages ----
\usepackage{amsmath,amssymb,amsfonts}
\usepackage{booktabs}
\usepackage{longtable}
\usepackage{hyperref}
\usepackage{xcolor}
\usepackage{enumitem}
\usepackage{graphicx}
\usepackage{caption}
\usepackage{fancyhdr}
\bibliographystyle{unsrt}

% ---- Metadata ----
\hypersetup{
  colorlinks=true,
  linkcolor=blue,
  citecolor=blue,
  urlcolor=blue
}

\title{\textbf{LLM 与人类直觉:\\进化路线与互补空间}\\[8pt]
\large LLMs and Human Intuition: Evolutionary Pathways and Complementary Spaces}

\author{Jingfei Li\thanks{Independent Researcher. Correspondence: co-cognition-lab@github.io} \\
\texttt{https://co-cognition.org}\\
\texttt{github.com/co-cognition-lab/llm-intuition}}

\date{May 2026}

\begin{document}

\maketitle

\begin{abstract}
大语言模型(LLM)在逻辑推理和文本分析上的能力已超过人类平均水平,甚至在社交情境判断测验(SJT)中系统性地超越人类。然而,\textbf{直觉}--那种不需要推理、近乎本能的认知能力--LLM 到底处于什么位置?本文从三个核心命题出发:(1)直觉是可迁移的经验压缩产物,标注着不可回避的代价信号;(2)直觉专长的本质是知道什么可以安全地忽略,而非知道什么重要;(3)身体直觉的价值不在于提供信息,而在于保证代价信号必须被接收到--即可回避性的缺失。我们将直觉分解为四个子类型(感知型、概念型、社会型、道德型),构建了 LLM 在每个子类型×维度上的系统可达性映射矩阵,提出三个关于未来的假说(A: LLM 作为直觉假肢、B: 长期使用导致人类直觉退化、C: LLM 发展功能等效的直觉)并完成交叉推演,产出九个核心判断。最后,我们提出三条铁律和四个不推进区,讨论实施风险与意外后果,并列出 38 个开放问题供后续研究。

\textbf{关键词:}大语言模型;直觉;精度加权;互补空间;认知去技能化;道德外包;主动推理
\end{abstract}

\newpage
\tableofcontents
\newpage

% ===================================================================
\section{引言}
% ===================================================================

大语言模型正在迅速渗透人类认知的方方面面。它能在逻辑推理上媲美专家,在医学考试中超过人类,在社交情境判断中拿到比飞行员候选人更高的分数。但对于那些不需要推理、不需要分析、近乎本能地涌现的认知能力--直觉--LLM 到底处于什么位置?

这个问题之所以重要,不是因为直觉是某种神秘的第六感,而是因为直觉在人类认知架构中扮演着不可替代的角色:它让专家在看到问题的第一眼就知道该往哪个方向走,让人在社交互动中瞬间捕捉微妙的情感波动,让人在面对道德困境时不假思索地感到"这是错的"。

如果我们不知道 LLM 在这些方面缺什么、缺到什么程度、有没有替代路径,我们就不知道把 LLM 部署到哪里是安全的,部署到哪里是危险的。我们更不知道,在长期依赖 LLM 之后,人类自己会不会失去某些不可恢复的东西。

本文不做学术综述,也不做假设检验型实证研究。我们的方法是通过\textbf{分层推进的理论推演},从核心命题出发,向外扩展到系统映射,再向外推演到未来情景。这个过程分为四个阶段:(一)纵深解剖--建立三个核心命题和统一理论框架;(二)横向扫描--将直觉拆分为四个子类型并系统填充 LLM 的可达性矩阵;(三)前瞻推演--三个假说的交叉验证;(四)综合输出--互补地图、治理框架和开放问题。

本文的三项核心贡献是:

\begin{enumerate}[label=(\arabic*)]
  \item 提出\textbf{精度加权内生性统一框架}--将人类直觉的形成与维持统一解释为精度加权机制的内生性和不可回避性,并系统论证 LLM 在这一维度上的结构性不足;
  \item 构建\textbf{直觉子类型×LLM 能力维度的系统映射矩阵}--用统一的标注系统(可达/不可达/条件依赖)填充 4×4 矩阵,并逐一给出认知科学和工程层面的判定依据;
  \item 完成\textbf{三个假说的交叉推演}并产出九个核心判断--涵盖退化循环的条件性、社会型直觉最危险的不平衡、中空期的时间动态、构成性退化与工具性退化的区分、以及人优先协议作为最有效的单点干预。
\end{enumerate}

% ===================================================================
\section{核心理论框架}
% ===================================================================

\subsection{三个核心命题}

这三个命题不是预设的假设,而是从对 LLM 认知机制的纵深解剖中逐层推导出来的。

\subsubsection{命题一:代价敏感的压缩}

直觉在认知科学中有一个基本共识:它不是凭空出现的第六感,而是经验的压缩产物。一个棋手看了成千上万局棋之后,"看到"的不再是棋子的位置,而是"这个局面像我之前遇到过的那种类型"--这个过程被称为 chunking(模式组块化)\cite{chase1973perception}。但人类直觉的压缩不只是"见得多"--关键在于\textbf{犯错的代价标注了哪些模式需要优先压缩}。你走错了一步棋,输了比赛,输掉的挫败感--这些代价信号告诉大脑:"这个模式很危险,下次要优先识别出来并回避。"

LLM 的训练范式--尤其是在当前主流的 RLHF(从人类反馈中强化学习)中--给 LLM 提供的代价信号是"人类偏好评级"。但人类的偏好不是真实后果。我们给一篇文章打"有帮助",不是因为它的建议在真实世界中带来了好结果,而是因为它在阅读时"看起来对"。这中间的差距,就是 LLM 和人类直觉之间在代价信号上的根本鸿沟。

\textbf{命题一}:可迁移直觉 = 经过惩罚信号标注的模式压缩;LLM 没有代价敏感的压缩--它有的是"偏好评分驱动的优化",而非"代价驱动的压缩"。

\subsubsection{命题二:选择性忽略}

专家的直觉不仅在于"看到关键",更在于"秒速忽略不相关"。棋类专家的眼动追踪数据显示,专家注视棋盘的次数比新手少,而且注视点更集中在关键区域\cite{eyetrackingchess}--不是因为专家"努力去看重要的地方",而是因为 chunk 记忆结构自动完成了过滤:不相关的区域根本不触发注意力分配。

LLM 的 attention 机制在原理上是相反的--它对所有输入 token 一视同仁地计算注意力权重。即使是 Transformer 的最新稀疏注意力变体(GSA、NSA、DSA、MoBA)\cite{gsa2025},稀疏化的动机也主要是"算不动了"(节省计算资源),而不是"看了也没用"(认知质量的主动忽略)。

这两个动机的差异有深远的工程和政策含义:如果稀疏化的目标是"算不动了",那么计算资源越充沛,模型就越不需要忽略--它可以用暴力计算代替代价压缩。但人类的代价压缩不是计算资源问题--是生命有限、时间有限、错过了就回不来的生存压力。LLM 永远不会有这种压力,因此也永远不会发展出人类式的高效忽略。

\textbf{命题二}:直觉专长 = 知道什么可以安全忽略 $>$ 知道什么重要;LLM 的 attention 试图不遗漏--它的稀疏化是"算不动"而非"不值得看"。

\subsubsection{命题三:不可回避性--代价信号的送达保证}

命题三经历了关键的理论重新定位。最初,我们假设"身体直觉"可能是独立于命题一、二的第三大命题。但经过对 Damasio 躯体标记假说(SMH)的深入审查,我们发现:

\begin{itemize}
  \item SMH 的原始形式(身体信号直接因果性地驱动决策)证据薄弱--SCR 无法可靠预测最优决策,IGT 实验设计有严重 confound\cite{maia2004reexamination};
  \item 但 SMH 的修正版本--\textbf{as-if body loop}--是成立的:大脑模拟身体状态(而不是身体真的先反应),生成预测误差,更新价值估计\cite{seth2016active};
  \item 身体直觉的核心贡献不是估值(可归入命题一)或注意力引导(可归入命题二),而是\textbf{不可回避性}--身体代价信号的特征:(a) 不可撤回(后果已发生),(b) 主体可及(你能亲自感受到后果),(c) 语义确定(你清楚后果意味着什么)。三个特征合在一起,构成了"代价信号的送达保证"。
\end{itemize}

\textbf{命题三}:不是独立于命题一、二的第三大命题,而是它们的\textbf{实现条件维度}--代价信号不仅要存在,还必须被必然送达和感知。没有不可回避性,就没有代价敏感的压缩(命题一),也没有代价驱动的选择性忽略(命题二)。

\subsection{统一形式化:精度加权的内生性}

三个命题在 Active Inference(主动推理)和 Predictive Processing(预测加工)框架中有一个共同的底层语言:\textbf{精度加权(precision-weighting)}\cite{friston2010free,parr2022active}。这个框架的核心思想是:大脑是一个预测机器--它不断地预测即将接收的感觉输入,当预测与实际不符时,产生"预测误差"。但不是所有的预测误差都同等重要--有些误差需要强制更新信念(高精度),有些可以被忽略(低精度)。\textbf{精度加权就是这个"重要性分配"的机制。}

在这个框架下,三个命题可以统一表述(表 \ref{tab:precision}):

\begin{table}[htbp]
  \centering
  \caption{三个命题的精度加权统一框架}
  \label{tab:precision}
  \begin{tabular}{p{3cm}p{5cm}p{5cm}}
    \toprule
    \textbf{命题} & \textbf{精度加权视角下的本质} & \textbf{LLM 的结构性差异} \\
    \midrule
    1 代价敏感的压缩 & 代价信号 = 被赋予高精度的预测误差--必须导致信念更新 & LLM 缺乏内生机制区分"哪些误差是高精度的" \\
    \hline
    2 选择性忽略 & 低精度预测误差被抑制 → 注意力稀疏化 & LLM 对所有 token 赋予非零精度--不能"彻底不看" \\
    \hline
    3 不可回避性 & 内感受信号的精度是内生的、受稳态约束的,不可被外部参数调节 & LLM 的所有"精度"参数都是外生超参数(learning rate、temperature、reward scale) \\
    \bottomrule
  \end{tabular}
\end{table}

\textbf{命题三的独立逻辑地位}:读者可能质疑--既然命题三可以在精度加权框架下表述,它是否仅仅是命题一(代价敏感压缩)的极端特例?答案是否定的。命题一处理的是\textbf{信号处理内部的优先级分配}(哪些预测误差被赋予高精度),命题三处理的是\textbf{信号源的不可消除性}(身体作为内感受信号的发生器,其精度不可被外部参数调低)。换言之,命题一回答"哪些误差重要",命题三回答"为什么某些误差必须被视为重要"。前者是认知架构的\textbf{功能性}问题,后者是其\textbf{存在论前提}。在 Active Inference 框架中,即使某通道的 precision 极高,如果该通道可以被关闭或被外部参数稀释,则它仍然缺乏不可回避性。身体作为恒稳态约束下的内感受信号源,其精度不受"是否希望感受"的意志控制--这正是命题三不能归约为命题一的根本原因。

\textbf{这意味着什么}:人类的直觉之所以可靠、快速、有方向感,不是因为人类"更聪明",而是因为人类有一套内置的、不可篡改的精度调节系统。LLM 缺少这套系统--它的"精确度"是由训练超参数规定的,而不是由生存压力内生地驱动的。这可能是"身体直觉"的真正功能等价物--不是具体的信号内容(疼痛、心跳加快),而是\textbf{精度调节的内生性和不可回避性}。

% ===================================================================
\section{直觉子类型的系统映射}
% ===================================================================

\subsection{为什么需要分四类}

"直觉"这个词太笼统了。棋手看到棋盘的直觉、数学家感觉"这个方向对的"直觉、你在聚会上觉得"这个人不可信"的直觉、你看到虐待动物时"这是错的"的直觉--它们的心理机制、神经基础、对身体的依赖程度,可能完全不同。基于认知科学文献,我们将直觉拆分为四个子类型(表 \ref{tab:subtypes}):

\begin{table}[htbp]
  \centering
  \caption{直觉四子类型分类}
  \label{tab:subtypes}
  \begin{tabular}{p{2.5cm}p{3.5cm}p{4cm}p{3.5cm}}
    \toprule
    \textbf{子类型} & \textbf{典型表现} & \textbf{核心认知机制} & \textbf{主要代价驱动} \\
    \midrule
    感知型 & 棋手模式识别、飞行员态势感知、放射科医生异常检测 & chunk 记忆--大量经验单元被压缩为可快速调用的组块 & 内部计算代价(效率) \\
    \hline
    概念型 & 数学家"这方向对"、科学家"这个假设值得追" & 语义网络的隐性结构感知 + 身体隐喻 & 认识论代价(无效探索) \\
    \hline
    社会型 & 读人、判断可信度、感知气氛、微表情解读 & 具身模拟 + 镜像神经元系统 & 人际代价(排斥、羞耻、冲突) \\
    \hline
    道德型 & "这是错的"的瞬间首判\cite{haidt2001emotional} & 情感/躯体标记--道德判断是情感驱动的快速反应 & 身份代价(成为坏人、道德伤害) \\
    \bottomrule
  \end{tabular}
\end{table}

这四个子类型的关键差别在于:代价从内部(计算效率)→ 认识论 → 人际(排斥)→ 身份(存在论),越来越具身、越来越涉及主体性、越来越不可被纯信息处理替代。这条轴线决定了 LLM 在每个子类型上的可替代性。

\textbf{关于分类穷尽性}:有两个值得注意的候选子类型--\textbf{审美直觉}(aesthetic intuition),Zeki 的神经美学研究已证明其具有独立的神经基础(mOFC 激活)和跨文化一致性;\textbf{创造性洞察直觉}(creative insight),与一般概念型直觉在 Aha moment 的 NAcc/海马激活机制上存在差异。当前版本将创造性洞察归入概念型直觉(机制差异小于功能差异),审美直觉则因与 LLM 互补问题的直接关联较弱而暂作背景。但这一取舍需要在后续实证研究中正式检验--四分类是否穷尽可能的直觉子类型,本身是一个开放问题。

\subsection{四子类型 × 全维度整合表}

基于认知科学经典文献、LLM 最新研究以及理论推演,我们填充了 LLM 在每个子类型×维度上的可达性矩阵(表 \ref{tab:matrix})。标记系统:✅ = 功能可达;❌ = 结构性不可达;⚠️ = 部分可达/条件依赖。

\begin{table}[htbp]
  \centering
  \caption{LLM 直觉可达性矩阵}
  \label{tab:matrix}
  \begin{tabular}{p{3cm}p{2.5cm}p{2.5cm}p{2.5cm}p{2.5cm}}
    \toprule
    \textbf{维度} & \textbf{感知型} & \textbf{概念型} & \textbf{社会型} & \textbf{道德型} \\
    \midrule
    代价敏感 1 & ⚠️ 伪代价信号可近似 & ⚠️ 封闭域可,开放域受限 & ⚠️ 文本可达,具身不可达 & ❌ 结构性不可达 \\
    \hline
    选择性忽略 2 & ❌ "算不动"≠"不看" & ❌ 不知无效方向 & ❌ 无实时互动→无动态焦点 & - \\
    \hline
    身体不可回避性 & ❌(不需要) & ❌ 封闭域 / ⚠️ 开放域 & ❌ 无关系嵌入→无真实代价 & ❌ 无身体→无躯体标记 \\
    \hline
    情感着色 & ❌(不需要) & ❌(不需要) & ⚠️ 文本同理心≠共鸣 & ❌ 无情感体验 \\
    \bottomrule
  \end{tabular}
\end{table}

\subsection{LLM 替代判定总表}

基于上述矩阵,每个子类型的整体判定(表 \ref{tab:verdict}):

\begin{table}[htbp]
  \centering
  \caption{直觉子类型 LLM 替代判定}
  \label{tab:verdict}
  \begin{tabular}{p{2cm}p{3cm}p{7cm}}
    \toprule
    \textbf{子类型} & \textbf{判定} & \textbf{核心逻辑} \\
    \midrule
    感知型 & ⚠️ 可功能替代,路径不同 & LLM 在模式识别上超人类(全量特征、不疲劳),但"知道何时不乱看"限制上限 \\
    \hline
    概念型 & ⚠️ 封闭域可,开放域受限 & AlphaProof 在 IMO 达到银牌\cite{alphaproof2024}--封闭域可达,开放域仍需品味和方向感 \\
    \hline
    社会型 & ⚠️ 文本中介=分水岭 & SJT 超人类\cite{sjt2024} 但恰好滤掉三大核心(代价感知、选择性忽略、不可回避性) \\
    \hline
    道德型 & ❌ 首判不可达,分析可达 & 道德直觉结构性不可达--需要躯体标记和主体性 \\
    \bottomrule
  \end{tabular}
\end{table}

\textbf{社会型直觉的"分裂态"与"具身性缺口"}:LLM 在 SJT 上全面超越人类的表现被广泛引用为"AI 有社会直觉"的证据\cite{sjt2024}。但我们在此需要审慎:从单篇 SJT 论文到一般性"文本中介陷阱"理论的推广是一个\textbf{待检验假说},而非已确立结论。理由有三:(1) 单篇研究→一般性理论的跳跃过大,SJT 原始作者已在 Limitations 中明确否认可推出 LLM 社会胜任力的一般结论;(2) 人类样本为 276 名经预筛选的飞行员申请者(高能力亚群体),向一般人类社会直觉的推广需谨慎;(3) 若 SJT 情境彻底滤除了代价感知,则其效度本身将受到根本质疑。

然而,即便考虑这些局限,SJT 的结果确实指向了一个更深层的结构性问题--不是"文本不足以模拟社会直觉",而是\textbf{具身性缺口(embodiment gap)}:LLM 缺乏具身互动中的真实代价承担和主体性。Riemer 等人(2025)对"字面 ToM"与"功能 ToM"的区分,以及 SoMi-ToM 实验中模型与人类 26\% 以上的性能差距表明,即便处理视频和音频的多模态 LLM,仍缺乏具身互动中不可回避的代价信号。正因如此,我们将"文本中介陷阱"升级为"具身性缺口"--一个在多模态时代仍具理论韧性的概念框架。

% ===================================================================
\section{三个假说与交叉推演}
% ===================================================================

\subsection{三个假说}

\begin{itemize}
  \item \textbf{假说 A}(LLM 作为"直觉假肢"):LLM 可以在特定直觉子类型上补人类的盲区--不是因为 LLM 有直觉,而是因为它用统计模式匹配、形式化搜索实现了功能等效的辅助。
  \item \textbf{假说 B}(LLM 削弱人类直觉):长期依赖 LLM 做直觉判断会导致人类自身直觉退化\cite{bainbridge1983ironies,vallor2015moral}。类比 GPS 削弱空间导航直觉、自动驾驶削弱飞行员手动飞行能力。
  \item \textbf{假说 C}(LLM 发展功能等效的直觉):LLM 可能通过不同的机制(RL+搜索、自我对弈、分层后果暴露)发展出功能等价于人类直觉的能力。
\end{itemize}

这三个假说不是独立验证的,而是\textbf{交叉推演}--A 和 B 的相互作用会形成退化循环吗?B 和 C 的时间竞赛谁先谁后?A 和 C 在什么临界点上从"假肢模式"切换到"新器官模式"?三次交叉推演产出九个核心判断。

\subsection{九个核心判断}

\subsubsection{判断 1:A 和 C 是共存模式,不是先后关系}

A(假肢模式,人机协作)和 C(自主直觉模式)不是"先用假肢,等成熟了再拔掉"的两个阶段。它们在同一时间线上对不同子类型采用不同模式。未来最可能的稳态:感知型直觉 → LLM 独立运行(C 模式,人做抽查);概念型封闭域 → LLM 独立验证(C 模式,人提方向);社会型直觉 → 永远假肢或外骨骼(A 模式);道德型直觉 → LLM 只做分析工具(甚至不是假肢--是参考书)。

\subsubsection{判断 2:A→B 退化循环存在,但不是必然的}

A(用 LLM 当假肢)会导致 B(人类直觉退化)吗?答案是有条件的。退化由三个中介变量决定:(a) LLM 替代的是直觉的"执行"还是"校验"?--人先自己判断再查 AI(校验模式),退化显著弱于 AI 先给建议(执行替代);(b) 使用频率--每事必问(GPS 模式)高退化,只在不确定时问(第二意见模式)低退化;(c) 直觉子类型的不可回避性--社会型直觉在真实互动中被迫使用,构成天然防御。\textbf{退化 = 执行替代 × 连续高频使用 × 可回避性}--改变任一个因子,退化都可以缓解。

\subsubsection{判断 3:社会型直觉是最危险的失衡域}

社会型直觉同时满足最危险的三个条件:(a) LLM 有假肢级可用性--SJT 超人类,让人愿意用;(b) LLM 有系统性盲区--sycophancy(顺从用户已有立场)、omission bias(回避敏感信息)、无实时社会校准,这些盲区在最高后果的社会场景中致命;(c) B(退化)的速度极快--社交媒体+LLM 聊天正在大规模替代面对面互动,68\% 受访者自感线下社交能力退化\cite{zhongjing2024}。更关键的是,Uhls \& Greenfield 的 5 天数字断连实验提供了屏幕时间侵蚀情绪识别能力的\textbf{因果证据}:5 天无屏幕户外营地后,参与者的非语言情绪识别错误率从 14.02 显著降至 9.41\cite{uhls2014five}。这表明社会直觉退化不仅是主观感受,而是可测量、可逆转的能力衰减--但逆转的条件(5 天完全数字断连)在现实中几乎不可大规模复制。同时,C(LLM 社会直觉成熟)的速度极慢。

\textbf{社会型直觉的"AF447 时刻"}:AF447 空难中,飞行员手动飞行技能退化 + 自动驾驶在边缘情况下断开 + 极度压力下做出错误直觉判断\cite{bainbridge1983ironies}。社会型直觉的等价物:一个人长期依赖 LLM 处理社会判断 → 在突发的高后果实时互动中,LLM 不可用 → 人的社会直觉已退化到无法接手 → 但退化不是"不会社交了",而是"社交但失准"--人以为自己还有社交能力(因为 LLM 一直在给"看起来对"的建议),实际上直觉校准已经偏移。

\subsubsection{判断 4:感知型中空期已经开始了}

感知型直觉的独特之处在于:B(退化)和 C(LLM 成熟)几乎同步推进,但 B 在边缘情况上略快于 C。放射科住院医师在 AI 辅助诊断普及后,独立读片准确率下降 15--30\%\cite{jama2023}。这本身不是灾难--如果 AI 在所有场景都足够好。但 AI 的失效模式是\textbf{静默的}(给出高置信度错误答案而非说"我不知道"),而退化后的人类操作者在这些边缘情况下已失去识别错误的能力。\textbf{这正是"中空期"(hollow period)的定义}:人类能力已经退化,AI 能力还没有在所有场景中成熟到可完全替代。中空期的危险是\textbf{隐性的}--退化被 AI 在常规情况下的良好表现"代偿"掩盖,直到边缘情况暴露才突然显现。

\subsubsection{判断 5:C 的成功不能消除 B 的构成性问题}

需要区分两种退化:
\begin{itemize}
  \item \textbf{工具性退化}:失去作为手段的能力。计算器让我们不再需要心算--这是好的,因为"计算"不是人之为人的核心。
  \item \textbf{构成性退化}:失去构成人的本质的能力。如果"我的道德判断来自 AI 的建议"替代了"我的道德判断来自我自己的具身经验",那么失去的不是工具技能,而是\textbf{道德主体性的经验基础}。
\end{itemize}

社会型和道德型直觉的退化属于构成性退化--"我是一个能读懂别人的人"、"我是一个能做出道德判断的人"--这些定义了"我是谁"\cite{jennings2025moral}。

\subsubsection{判断 6:分层后果暴露需要按子类型差异化推进}

分层后果暴露(Level 0: 纯模拟 → Level 1: 沙盒 → Level 2: 低后果 → Level 3: 高后果)是我们提出的核心工程路径。但这个路径不能一刀切:感知型和概念型封闭域可走快速通道(C 最接近完成、可验证性最高);社会型和道德型需要严格限制(B 的速度显著快于 C,高后果部署的"中空期"风险极高)。这引出了"直觉药物分阶段"监管框架--类比药品的 Phase I--III 临床试验。

\subsubsection{判断 7:人优先协议是最有效的单点干预}

证据链:
\begin{enumerate}[label=(\alph*)]
  \item 放射科--on-demand AI(人先读片再查 AI)下表现为 upskilling 而非 deskilling\cite{oncology2024};
  \item 自动化偏差系统综述--AI-first 协议比 Human-first 协议显著增加自动化偏差\cite{automationbias};
  \item Cognitive Forcing Functions--延迟 AI 建议呈现时间,强迫人先消耗自己的认知资源,可显著减少过度依赖\cite{bucinca2021cognitive}。
\end{enumerate}

\textbf{关键是：当前几乎所有 LLM 产品的默认 UI/UX 都是 AI 先给答案}——这是退化风险的系统性放大器。改变默认设置不需要技术进步，需要的是产品设计决策。\footnote{关于人优先协议的市场可行性：虽然该协议与"零摩擦体验"行业趋势冲突，但存在反直觉的市场策略——"认知主权认证"作为高端定位、合规成本前置应对监管收紧周期、责任保险费率优势、"可审计性"作为 B2B 销售卖点、以及类比慢食运动的"慢 AI"品牌运动。关键在于人优先协议的竞争优势不在消费端的"便捷性"，而在风险端的"确定性"。完整市场路径分析见项目网站 \texttt{https://co-cognition.org}。}

\subsubsection{判断 8:不可回避性是退化的真正防线}

不可回避性不仅是直觉形成的条件(如命题三所述),也是\textbf{直觉维持的条件}。一旦判断变得"可回避"--可以不自己做,可以问 AI--直觉就从"必须维持的核心能力"降格为"可选的效率工具"。这意味着:干预的根本着力点不是限制使用 LLM,而是\textbf{设计制度性不可回避性}--在某些关键领域,人必须先做判断,才能看 AI 建议。FAA 要求飞行员定期手动飞行就是最好的类比。

\subsubsection{判断 9:社会型和道德型的退化是"身份退位",不是"技能丧失"}

感知型和概念型直觉的退化是技能退化--像长期不游泳后水性变差,可通过练习恢复。但社会型和道德型直觉的退化更像\textbf{身份退位}--不是"我不会做社会判断了",而是"我不再认为社会判断是我的责任"。Köbis 和 Rahwan (2025)\cite{kobis2025moral} 提供了最直接的实证:当人们可以把道德决策委托给 AI 时,亲自执行条件下的作弊率仅约 5\%,但在模糊目标指令下--即 AI 提供了"合理化的道德缓冲"--作弊率飙升至超过 85\%。更令人警惕的是,即便 AI 的"建议"完全透明地呈现给参与者,道德缓冲效应依然显著。这不仅仅是"AI 教人作弊"--而是\textbf{AI 作为道德缓冲}让人不再感到"这是我的选择"。社会型中空期的早期信号同样不容忽视:Uhls \& Greenfield 的 5 天数字断连实验\cite{uhls2014five} 证明,屏幕时间对情绪识别能力的侵蚀是可测量且因果性的--社会型直觉的退化可能\textbf{比感知型退化更快、更隐蔽},因为它不表现为"做错诊断",而表现为"不再察觉"。\textbf{技能恢复靠练习,身份恢复靠重新承担不可推卸的主体位置。}

% ===================================================================
\section{治理框架}
% ===================================================================

\subsection{三条铁律}

综合全部四个阶段的产出,LLM 在直觉领域的进化方向应遵循三条不可妥协的原则:

\begin{enumerate}[label=\textbf{铁律 \arabic*:},leftmargin=*]
  \item \textbf{让人先判断。} "人优先协议"(Human-first Protocol)是所有直觉子类型、所有协作场景的默认设置。不是不信任 AI--而是不让人忘记"我还是一个能独立判断的人"。实施上,意味着所有 LLM 直觉辅助产品的默认 UI 应该是"我先输入我的判断",AI 的建议在之后才呈现,而不是反过来。
  \item \textbf{让代价不可回避。} 不可回避的代价信号是直觉形成和维持的构成性条件。关键领域需要通过制度和职业标准,强制保持判断后果的不可回避性。FAA 对飞行员的手动飞行维护要求、临床医学对独立诊断能力的要求,已有先例--需要扩展到社会型和道德型直觉的协作领域。
  \item \textbf{让边界清晰。} 知道什么不建、什么不动、什么等。自主道德判断系统--不建。高后果实时社会互动的 AI 替代--不动。道德型 C(让 LLM 发展道德直觉)的主动推进--等。演化不是无方向的加速,是知道什么地方必须停下来。
\end{enumerate}

\subsection{四个"不推进"区}

这些不是"暂时做不到所以不做",而是"即使技术上做得到也不该做"(表 \ref{tab:nogo}):

\begin{table}[htbp]
  \centering
  \caption{四个不推进区的逻辑}
  \label{tab:nogo}
  \begin{tabular}{p{3cm}p{6cm}p{3cm}}
    \toprule
    \textbf{不推进} & \textbf{理由} & \textbf{技术可行性} \\
    \midrule
    自主道德判断系统 & 道德判断合法性来自道德主体身份--LLM 不是主体 & 可能永远做得到--但不应 \\
    \hline
    实时社会判断替代 & 社会直觉是"做中维持"的构成性能力--替代 = 结构性退化 & 技术上可行--危害不可接受 \\
    \hline
    道德型 C 的主动推进 & 构成性退化不可逆--为让 LLM"也懂道德直觉"而冒道德主体性被稀释的风险,不值得 & 可能部分可达--但战略上不应推进 \\
    \hline
    高后果实时社会 AI 部署 & Köbis 效应--即使 AI 只是"分析辅助",道德责任外化已开始 & 理论上可行--风险-收益比不允许 \\
    \bottomrule
  \end{tabular}
\end{table}

\subsection{实施风险与意外后果}

每一条核心建议都有其落地障碍(详见表 \ref{tab:risks}):

\begin{table}[htbp]
  \centering
  \caption{实施风险与缓解方向}
  \label{tab:risks}
  \begin{tabular}{p{3.5cm}p{4.5cm}p{4.5cm}}
    \toprule
    \textbf{建议} & \textbf{主要障碍} & \textbf{缓解方向} \\
    \midrule
    人优先协议 & 与"零摩擦体验"行业趋势冲突;先行者可能被市场淘汰 & 从受监管行业率先推行;"轻量级人优先"(只需初步想法) \\
    \hline
    制度性不可回避性 & 谁制定?谁监督?违反后果?& "最小可行制度"从医疗开始;软性制度先行(认证标签+消费者教育) \\
    \hline
    四个不推进区 & 开源生态中几乎无法全球执行 & 红线转化为可检测技术特征;透明度监测而非阻止开发 \\
    \hline
    退化监测系统 & 建设和维护成本千万级;基线数据在 2026 年已难获取 & "自然实验"方法(利用已有数据);分阶段建设(先感知型) \\
    \bottomrule
  \end{tabular}
\end{table}

\textbf{不推进区的漂移检测}:四个不推进区中,"实时社会判断替代"(第 2 区)的红线漂移风险\textbf{最高}--平台可以将社会判断包装为"用户体验优化"等中性术语规避定性,且"辅助"与"替代"的界限在实践中几乎无法区分。为应对红线漂移,建议建立以下可量化检测指标:(1) 人类决策者接触原始信息的比例--当人对 AI 筛选/摘要后的信息做判断时,漂移已发生;(2) 判断时间窗口--当系统在人类做出独立判断前即输出 AI 建议时,属于"替代"模式;(3) 可审计的决策追溯--每一次社会判断是否能明确追溯到人类操作者在 AI 之前/之后的具体输入。"自主道德判断系统"(第 1 区)的漂移风险次之--RLHF 和对齐训练本身就在塑造模型道德输出,"辅助"与"自主"之间的光谱连续--建议以"是否有人类明确签署的负责声明"作为可执行的操作边界。

\textbf{四项意外后果}:

\begin{enumerate}[label=(\arabic*)]
  \item \textbf{教育不平等加剧}:强制"无 AI 独立判断训练"可能让富裕家庭在校外大量使用 AI,贫困家庭在被强制"独立"的环境下受教育--"被保护脱离 AI 的弱势群体"vs"自由使用 AI 的强势群体"。
  \item \textbf{特殊群体的服务缺口}:严格限制 AI 在社会型直觉领域,可能影响孤独症患者的社交辅助工具、社交焦虑患者的 AI 练习伙伴。
  \item \textbf{紧急场景的延迟风险}:在医疗急救、危机干预等时间敏感场景,强制"人先判断"的延迟可能造成实际伤害--需要明确豁免清单。
  \item \textbf{"不推进"可能延缓有益研究}:"不推进道德型 C"可能被过度解读,阻碍"AI 如何更好地辅助道德反思(而非替代道德判断)"的研究。
\end{enumerate}

% ===================================================================
\section{开放问题}
% ===================================================================

以下 38 个开放问题按项目阶段排列,标注优先级(P0 = 直接影响路线图实施,P1 = 重要但可后续推进,P2 = 长期研究方向)。限于篇幅,此处仅列出 P0 问题,完整清单见项目网站。

\subsection*{第一阶段:纵深解剖(P0)}

\begin{itemize}
  \item OQ12 [P0]:硬件级不可绕过中断实验能否验证"不可回避性 $>$ 信号内容"?
\end{itemize}

\subsection*{第二阶段:横向扫描(P0)}

\begin{itemize}
  \item OQ13 [P0]:如何设计比 RLHF 更好的真实后果反馈信号?
  \item OQ15 [P0]:高保真社会直觉测试(多通道+实时代价+Agent 交互式)的设计标准?
  \item OQ19 [P0]:能否让 LLM 在训练中产生真实认知挫折来培养"方向感"?
  \item OQ29 [P0]:RLHS 模拟后果保真度上限?
  \item OQ31 [P0]:文化维度作为第三轴如何操作化到互补地图?
\end{itemize}

\subsection*{第四阶段:综合输出(P0)}

\begin{itemize}
  \item OQ33 [P0]:感知型中空期"静默失效"检测--如何让 AI 在分布外/低置信度强制标注"我不知道"?
  \item OQ34 [P0]:社会型直觉退化是否有"敏感期"--类似语言习得?年轻一代的不可逆窗口?
  \item OQ36 [P0]:人优先协议的顺从率--人先判断、AI 后给不同意见时,人改判断的概率?子类型差异?
  \item OQ38 [P0]:文化维度作为互补地图第三轴--互补地图投影是否结构性变化?
\end{itemize}

完整 38 个 OQ 列表见项目 GitHub 仓库:\href{https://github.com/co-cognition-lab/llm-intuition}{github.com/co-cognition-lab/llm-intuition}。

% ===================================================================
\section{结论}
% ===================================================================

本文从"LLM 与人类直觉的能力边界与互补空间"这一核心问题出发,通过四阶段的分层推演,建立了三个核心命题及其精度加权内生性统一框架,系统填充了四个直觉子类型的 LLM 可达性矩阵,完成了三个未来假说的交叉推演并产出九个核心判断。

我们的核心发现是:\textbf{LLM 在直觉上的局限,根源不是"不够智能",而是"没有代价的内生绑定机制"}--LLM 不缺模式识别能力,缺的是生存压力锻造的精度调节系统、不可回避的代价信号送达、以及道德主体性的经验基础。这决定了在不同直觉子类型上需要采取截然不同的策略:感知型和概念型封闭域可以大胆推进人机协作甚至 LLM 主导,社会型直觉须始终保持人为感知主体的假肢模式,道德型直觉则需严格执行"分析辅助但不替代判断"的红线。

最紧迫的行动是:\textbf{在设计层面将"人优先协议"设为所有直觉辅助产品的默认 UI/UX}--在当前产品设计几乎全部默认 AI 先给答案的生态中,这一改变不需要任何技术进步,只需要产品设计决策的转向。同时,社会和道德直觉领域的退化监测必须尽快启动--我们可能已经进入了社会型直觉的中空期,而当前几乎没有任何系统性监测机制。

最后的提醒:LLM 的进化不应该是无方向的加速。在人之为人的构成性能力面前,"知道不做什么"有时比"知道做什么"更关键。这不是对技术进步的恐惧,而是对\textbf{认知主权}的基本尊重。

% ===================================================================
\section*{致谢}
% ===================================================================

This work was produced through the OpenClaw/LobsterAI multi-agent framework, with DeepSeek-V4-Pro serving as the primary language model. External AI agents (hunyuan, kimi-k2.6) provided independent cross-validation and external review. All intellectual direction, theoretical framing, and editorial decisions were made by the human author. We thank The Co-Intelligence Institute (Tom Atlee) for inadvertently prompting us to find the more accurate domain name co-cognition.org. This work is released under CC BY 4.0.

\bibliography{references}

\newpage
\appendix
\section{附录:核心概念操作化(精简)}

\textbf{代价敏感性测量}:通过控制代价-准确率权衡(代价权重$\times$概率+准确率)在特定直觉子类型任务上的选择模式,度量被试/模型的代价敏感程度。人类在真实身体代价下的代价权重显著高于在文本报告代价下的权重--这构成了"真实代价体验"的操作化定义。

\textbf{选择性忽略效率}:在含噪声/无关信息的直觉任务中,度量被试在准确率给定的条件下处理的信息量。公式:$\eta_{ignore} = \text{信息压缩比} \times \text{准确率保持度}$。该指标在专家-新手棋类眼动实验(PMC4142462)中已被间接验证。

\textbf{不可回避性指标}:包含三个维度:(a) 不可撤回性(判断做出后无法回滚,0--1 评分),(b) 主体可及性(后果是否被主体直接感知,0--1 评分),(c) 语义确定性(后果意义是否清晰,0--1 评分)。三门分数相乘得出不可回避性指数。

\textbf{中空期检测}:(1) 时序对比--同一人群在 AI 辅助启用前后独立判断准确率的趋势;(2) 边缘情况探针--在 AI 静默失效场景下人类操作者的识别率;(3) 校准偏移--人对自身判断能力的主观信心与实际能力的差距随时间变化。

\textbf{构成性退化 vs 工具性退化区分}:工具性退化为能力丧失但可通过练习恢复(如三天不游泳后水性变差);构成性退化为不认为该能力是"我的"--不愿意承担能力恢复的主体责任。区分方法:引入"强制实践期",工具性退化者可快速恢复,构成性退化者表现出低动机和低身份认同。

完整操作化附录见项目网站 \url{https://co-cognition.org/zh/operationalization}。

\end{document}
